\documentclass[11pt]{article}
\usepackage[margin=1in]{geometry}
\usepackage[T1]{fontenc}
\usepackage{lmodern}
\usepackage{amsmath,amssymb}
\usepackage{booktabs}
\usepackage{siunitx}
\usepackage{graphicx}
\usepackage{xcolor}
\usepackage[numbers,sort&compress]{natbib}
\usepackage{hyperref}
\usepackage{cleveref}
\hypersetup{colorlinks=true,linkcolor=blue!50!black,citecolor=blue!50!black,urlcolor=blue!50!black}
\graphicspath{{../figures/}}
\newcommand{\fig}[1]{\IfFileExists{../figures/#1.pdf}{\includegraphics[width=0.88\linewidth]{#1.pdf}}%
  {\fbox{\parbox{0.9\linewidth}{\centering figure #1 not generated yet}}}}

\renewcommand{\topfraction}{0.9}
\renewcommand{\bottomfraction}{0.8}
\renewcommand{\textfraction}{0.07}
\renewcommand{\floatpagefraction}{0.85}
\setcounter{topnumber}{3}
\setcounter{totalnumber}{4}

% GENERATED by figures/make_figures.py -- do not edit by hand.
% benchmark source: results/dev.json; validation: results/validation.json
\newif\ifprovisional
\provisionaltrue
\newcommand{\benchDate}{2026-10-09}
\newcommand{\benchSource}{dev.json}
\newcommand{\bestFpThirtyTwoGBs}{221}
\newcommand{\bestFpThirtyTwoMLUPS}{3{,}076}
\newcommand{\bestFpThirtyTwoName}{Warp fused}
\newcommand{\bestFpThirtyTwoPct}{67}
\newcommand{\bestGBs}{205}
\newcommand{\bestMLUPS}{5{,}693}
\newcommand{\bestName}{CUDA C fused fp16}
\newcommand{\bestPctCopy}{62}
\newcommand{\bestPctSpec}{46}
\newcommand{\bestSixteenMLUPS}{5{,}693}
\newcommand{\bestSixteenName}{CUDA C fused fp16}
\newcommand{\bestSixteenPct}{62}
\newcommand{\bestThreeDMLUPS}{1{,}692}
\newcommand{\bestThreeDName}{Numba-CUDA fused}
\newcommand{\bestThreeDPct}{78}
\newcommand{\bestThreeDSixteenMLUPS}{2{,}304}
\newcommand{\bestThreeDSixteenPct}{53}
\newcommand{\ceilFpThirtyTwo}{4{,}587}
\newcommand{\ceilSixteen}{9{,}174}
\newcommand{\ceilThreeD}{2{,}173}
\newcommand{\compileRollMLUPS}{2{,}074}
\newcommand{\compileRollPct}{45}
\newcommand{\copyBW}{330}
\newcommand{\cudaFusedMLUPS}{3{,}029}
\newcommand{\cudaFusedPct}{66}
\newcommand{\eagerRollMLUPS}{108}
\newcommand{\eagerRollPct}{2}
\newcommand{\eqCount}{15}
\newcommand{\eqMaxBf}{2.1e-04}
\newcommand{\eqMaxFpSixteen}{3.6e-05}
\newcommand{\eqMaxFpThirtyTwo}{1.7e-06}
\newcommand{\ghiaFourLinf}{0.0047}
\newcommand{\ghiaFourLtwo}{0.0024}
\newcommand{\ghiaFourN}{256}
\newcommand{\ghiaOneBfLinf}{0.0285}
\newcommand{\ghiaOneBfLtwo}{0.0161}
\newcommand{\ghiaOneBfN}{128}
\newcommand{\ghiaOneHalLinf}{0.0050}
\newcommand{\ghiaOneHalLtwo}{0.0025}
\newcommand{\ghiaOneHalN}{128}
\newcommand{\ghiaOneLinf}{0.0059}
\newcommand{\ghiaOneLtwo}{0.0029}
\newcommand{\ghiaOneN}{128}
\newcommand{\ghiaOneOrigLinf}{0.0718}
\newcommand{\ghiaOneOrigLtwo}{0.0250}
\newcommand{\ghiaOneOrigN}{128}
\newcommand{\ghiaOneSixtyfourLinf}{0.0060}
\newcommand{\ghiaOneSixtyfourLtwo}{0.0028}
\newcommand{\ghiaOneSixtyfourN}{64}
\newcommand{\ghiaOneTwofiftysixLinf}{0.0072}
\newcommand{\ghiaOneTwofiftysixLtwo}{0.0033}
\newcommand{\ghiaOneTwofiftysixN}{256}
\newcommand{\ghiaThouAALinf}{0.0080}
\newcommand{\ghiaThouAALtwo}{0.0040}
\newcommand{\ghiaThouAAN}{256}
\newcommand{\ghiaThouBfLinf}{0.0653}
\newcommand{\ghiaThouBfLtwo}{0.0413}
\newcommand{\ghiaThouBfN}{256}
\newcommand{\ghiaThouCudaAAHalLinf}{0.0129}
\newcommand{\ghiaThouCudaAAHalLtwo}{0.0072}
\newcommand{\ghiaThouCudaAAHalN}{256}
\newcommand{\ghiaThouHalLinf}{0.0141}
\newcommand{\ghiaThouHalLtwo}{0.0079}
\newcommand{\ghiaThouHalN}{256}
\newcommand{\ghiaThouLinf}{0.0080}
\newcommand{\ghiaThouLtwo}{0.0040}
\newcommand{\ghiaThouMRTLinf}{0.0082}
\newcommand{\ghiaThouMRTLtwo}{0.0040}
\newcommand{\ghiaThouMRTN}{256}
\newcommand{\ghiaThouN}{256}
\newcommand{\ghiaThouTRTLinf}{0.0074}
\newcommand{\ghiaThouTRTLtwo}{0.0040}
\newcommand{\ghiaThouTRTN}{256}
\newcommand{\gpuName}{NVIDIA GeForce RTX 3070 Laptop GPU}
\newcommand{\karmanCd}{3.312}
\newcommand{\karmanCl}{1.005}
\newcommand{\karmanD}{40}
\newcommand{\karmanNx}{880}
\newcommand{\karmanNy}{164}
\newcommand{\karmanSt}{0.2996}
\newcommand{\karmanTau}{0.560}
\newcommand{\ladderRows}{D2Q9 & NumPy (fp64) & $512\times512$ & 11 & 144 & 2 & 0 & 0 \\
D2Q9 & torch eager, roll & $512\times512$ & 108 & 72 & 8 & 2 & 2 \\
D2Q9 & torch eager, gather & $512\times512$ & 117 & 72 & 8 & 3 & 2 \\
D2Q9 & original code (torch.compile) & $2048\times2048$ & 271 & 72 & 19 & 6 & 4 \\
D2Q9 & torch.compile, roll & $2048\times2048$ & 2{,}074 & 72 & 149 & 45 & 33 \\
D2Q9 & torch.compile, gather & $2048\times2048$ & 636 & 72 & 46 & 14 & 10 \\
D2Q9 & Numba-CUDA fused & $2048\times2048$ & 2{,}000 & 72 & 144 & 44 & 32 \\
D2Q9 & Warp fused & $2048\times2048$ & 3{,}076 & 72 & 221 & 67 & 49 \\
D2Q9 & Triton fused & $2048\times2048$ & 2{,}109 & 72 & 152 & 46 & 34 \\
D2Q9 & CUDA C fused & $2048\times2048$ & 3{,}029 & 72 & 218 & 66 & 49 \\
D2Q9 & Triton fused AA & $2048\times2048$ & 1{,}863 & 72 & 134 & 41 & 30 \\
D2Q9 & CUDA C fused AA & $2048\times2048$ & 2{,}081 & 72 & 150 & 45 & 33 \\
D2Q9 & Triton fused fp16 & $2048\times2048$ & 5{,}314 & 36 & 191 & 58 & 43 \\
D2Q9 & Triton fused bf16 & $2048\times2048$ & 3{,}787 & 36 & 136 & 41 & 30 \\
D2Q9 & CUDA C fused fp16 & $2048\times2048$ & 5{,}693 & 36 & 205 & 62 & 46 \\
D2Q9 & Triton AA fp16 & $2048\times2048$ & 5{,}681 & 36 & 205 & 62 & 46 \\
D2Q9 & CUDA C AA fp16 & $2048\times2048$ & 4{,}667 & 36 & 168 & 51 & 38 \\
\midrule
D3Q19 & NumPy (fp64) & $32\times32\times32$ & 5 & 304 & 1 & 0 & 0 \\
D3Q19 & torch.compile, roll & $64\times64\times64$ & 839 & 152 & 128 & 39 & 28 \\
D3Q19 & Numba-CUDA fused & $128\times128\times128$ & 1{,}692 & 152 & 257 & 78 & 57 \\
D3Q19 & Warp fused & $128\times128\times128$ & 1{,}355 & 152 & 206 & 62 & 46 \\
D3Q19 & Triton fused & $128\times128\times128$ & 1{,}022 & 152 & 155 & 47 & 35 \\
D3Q19 & CUDA C fused & $128\times128\times128$ & 1{,}007 & 152 & 153 & 46 & 34 \\
D3Q19 & Triton fused AA & $128\times128\times128$ & 1{,}025 & 152 & 156 & 47 & 35 \\
D3Q19 & CUDA C fused AA & $128\times128\times128$ & 899 & 152 & 137 & 41 & 30 \\
D3Q19 & Triton fused fp16 & $128\times128\times128$ & 2{,}036 & 76 & 155 & 47 & 35 \\
D3Q19 & CUDA C fused fp16 & $128\times128\times128$ & 2{,}304 & 76 & 175 & 53 & 39 \\
D3Q19 & CUDA C AA fp16 & $128\times128\times128$ & 1{,}389 & 76 & 106 & 32 & 24 \\}
\newcommand{\numbaFusedMLUPS}{2{,}000}
\newcommand{\numbaFusedPct}{44}
\newcommand{\numpyRefMLUPS}{11}
\newcommand{\numpyRefPct}{0}
\newcommand{\origMLUPS}{271}
\newcommand{\origPct}{5.9}
\newcommand{\poisBgkFinest}{1.7e-04}
\newcommand{\poisFpSixtyFourHThirtyTwo}{7.0e-04}
\newcommand{\poisFpThirtyTwo}{1.1e-03}
\newcommand{\poisOrder}{2.000}
\newcommand{\poisTrtMax}{4e-12}
\newcommand{\shearMaxErr}{0.17}
\newcommand{\shearMidErr}{0.09}
\newcommand{\shearN}{128}
\newcommand{\shearTauMax}{1.95}
\newcommand{\shearTauMin}{0.52}
\newcommand{\specBW}{448}
\newcommand{\speedupBest}{21.0}
\newcommand{\speedupFpThirtyTwo}{11.4}
\newcommand{\tgvMaxErr}{0.11}
\newcommand{\tgvN}{64}
\newcommand{\tritonFusedMLUPS}{2{,}109}
\newcommand{\tritonFusedPct}{46}
\newcommand{\tritonMRTMLUPS}{1{,}513}
\newcommand{\tritonMRTPct}{33}
\newcommand{\tritonTRTMLUPS}{2{,}129}
\newcommand{\tritonTRTPct}{46}
\newcommand{\warpFusedMLUPS}{3{,}076}
\newcommand{\warpFusedPct}{67}

% fallbacks for macros that a partial results JSON may not define
\providecommand{\origMLUPS}{--}\providecommand{\origPct}{--}
\providecommand{\speedupFpThirtyTwo}{--}\providecommand{\speedupBest}{--}
\providecommand{\bestThreeDMLUPS}{--}\providecommand{\bestThreeDName}{--}
\providecommand{\bestThreeDPct}{--}\providecommand{\bestThreeDSixteenMLUPS}{--}
\providecommand{\bestThreeDSixteenPct}{--}
\providecommand{\tritonTRTPct}{--}\providecommand{\tritonMRTPct}{--}
\providecommand{\tritonTRTMLUPS}{--}\providecommand{\tritonMRTMLUPS}{--}
\providecommand{\compileRollMLUPS}{--}\providecommand{\compileRollPct}{--}
\providecommand{\eagerRollMLUPS}{--}\providecommand{\numpyRefMLUPS}{--}
\providecommand{\tritonFusedPct}{--}\providecommand{\warpFusedPct}{--}
\providecommand{\numbaFusedPct}{--}\providecommand{\cudaFusedPct}{--}
\providecommand{\tritonFusedMLUPS}{--}\providecommand{\warpFusedMLUPS}{--}
\providecommand{\numbaFusedMLUPS}{--}\providecommand{\cudaFusedMLUPS}{--}
\providecommand{\ghiaOneLtwo}{--}\providecommand{\ghiaOneLinf}{--}
\providecommand{\ghiaFourLtwo}{--}\providecommand{\ghiaFourLinf}{--}
\providecommand{\ghiaThouLtwo}{--}\providecommand{\ghiaThouLinf}{--}
\providecommand{\ghiaThouMRTLtwo}{--}\providecommand{\ghiaThouMRTLinf}{--}
\providecommand{\ghiaThouTRTLtwo}{--}\providecommand{\ghiaThouTRTLinf}{--}
\providecommand{\ghiaOneHalLtwo}{--}\providecommand{\ghiaOneHalLinf}{--}
\providecommand{\ghiaOneBfLtwo}{--}\providecommand{\ghiaOneBfLinf}{--}
\providecommand{\ghiaThouHalLtwo}{--}\providecommand{\ghiaThouHalLinf}{--}
\providecommand{\ghiaThouBfLtwo}{--}\providecommand{\ghiaThouBfLinf}{--}
\providecommand{\ghiaThouAALtwo}{--}\providecommand{\ghiaThouAALinf}{--}
\providecommand{\ghiaThouCudaAAHalLtwo}{--}\providecommand{\ghiaThouCudaAAHalLinf}{--}
\providecommand{\ghiaOneOrigLtwo}{--}\providecommand{\ghiaOneOrigLinf}{--}
\providecommand{\ghiaOneSixtyfourLtwo}{--}\providecommand{\ghiaOneSixtyfourLinf}{--}
\providecommand{\ghiaOneTwofiftysixLtwo}{--}\providecommand{\ghiaOneTwofiftysixLinf}{--}
\providecommand{\shearMaxErr}{--}\providecommand{\shearMidErr}{--}\providecommand{\shearN}{--}
\providecommand{\shearTauMin}{--}\providecommand{\shearTauMax}{--}
\providecommand{\poisOrder}{--}\providecommand{\poisBgkFinest}{--}
\providecommand{\poisTrtMax}{--}\providecommand{\poisFpThirtyTwo}{--}
\providecommand{\poisFpSixtyFourHThirtyTwo}{--}
\providecommand{\karmanSt}{--}\providecommand{\karmanCd}{--}\providecommand{\karmanCl}{--}
\providecommand{\karmanD}{--}\providecommand{\karmanTau}{--}
\providecommand{\karmanNx}{--}\providecommand{\karmanNy}{--}
\providecommand{\tgvMaxErr}{--}\providecommand{\tgvN}{--}
\providecommand{\eqMaxFpThirtyTwo}{--}\providecommand{\eqMaxFpSixteen}{--}
\providecommand{\eqMaxBf}{--}\providecommand{\eqCount}{--}
\providecommand{\bestSixteenMLUPS}{--}\providecommand{\bestSixteenName}{--}
\providecommand{\bestSixteenPct}{--}

\title{Reaching the memory roofline with Python:\\
a fused-kernel lattice-Boltzmann solver on a laptop GPU}
\author{Walid Hemati\\ \small Independent researcher}
\date{}

\begin{document}
\maketitle

\ifprovisional
\begin{center}
\fcolorbox{red}{red!8}{\parbox{0.92\linewidth}{\centering\bfseries
PROVISIONAL NUMBERS\\[2pt]
\normalfont\small Performance figures in this draft come from a development benchmark run
(\texttt{results/\benchSource}) on a GPU shared with other jobs. They will be replaced by the
serial \texttt{bench/run.py --full} run on an idle GPU.}}
\end{center}
\fi

\begin{abstract}
Lattice-Boltzmann (LB) solvers are bandwidth-bound: one D2Q9 update in single precision
moves at least \SI{72}{B} through DRAM, so the attainable speed is fixed by the memory
system, not by the arithmetic. We ask whether a solver written and driven entirely from
Python can sit on that roofline. Starting from an existing PyTorch cavity code, we build
\texttt{pylbm}, a validated D2Q9/D3Q19 solver with BGK, TRT and MRT collisions, and climb an
eight-rung optimization ladder from NumPy through eager and compiled PyTorch to fused
stream--collide kernels in Triton, NVIDIA Warp, Numba-CUDA and hand-written CUDA C. On an
RTX~3070 Laptop GPU the best fp32 D2Q9 kernel (\bestFpThirtyTwoName) runs at
\bestFpThirtyTwoMLUPS{} MLUPS, \bestFpThirtyTwoPct\% of the measured device copy bandwidth
of \SI{\copyBW}{GB/s}; storing shifted populations $f_i-w_i$ in 16 bits halves the traffic and
reaches \bestSixteenMLUPS{} MLUPS. D3Q19 reaches \bestThreeDMLUPS{} MLUPS
(\bestThreeDPct\% of copy bandwidth). The Triton, Warp and Numba kernels are generated from a
single symbolic description, so all four toolchains execute identical arithmetic. The
solver is checked against Ghia et al.\ (1982), shear-wave decay, Poiseuille flow, the
Sch\"afer--Turek cylinder benchmark and 3D Taylor--Green decay; along the way we document
three boundary-condition defects of the original code.
\end{abstract}

\section{Introduction}
The lattice-Boltzmann method (LBM)~\citep{kruger2017,qian1992} is a standard tool for
weakly compressible flows. Its update is local and regular, which makes it a natural fit for
GPUs~\citep{tolke2008,obrecht2011,mawson2014}, and its arithmetic intensity is low: a D2Q9
node reads nine populations and writes nine, about two floating-point operations per byte
in fp32. On every current GPU this places LBM firmly under the memory roof of the roofline
model~\citep{williams2009}. The relevant figure of merit is therefore the fraction of the
attainable memory bandwidth a code sustains, not its FLOP rate.

Python is the language in which many scientists prototype, but it is rarely the language of
record for production LB codes, which tend to be C++/CUDA~\citep{bauer2021walberla} or use
Python only to generate them~\citep{bauer2021lbmpy}. Python LB frameworks built on tensor
libraries~\citep{bedrunka2021,ataei2024} trade some performance for flexibility. This paper
starts from a small PyTorch lid-driven-cavity script (the earlier version of the present
repository) that ran at a few hundred MLUPS and asks how far the same machine can be pushed
without leaving Python as the driving language.

Contributions:
\begin{itemize}
\item A restructured, tested solver (\texttt{pylbm}) with a consistent axis convention,
half-way bounce-back with Ladd's moving-wall correction~\citep{ladd1994}, BGK, TRT and MRT
collision, Guo forcing~\citep{guo2002}, a pure NumPy float64 reference implementation that
acts as the correctness oracle, and a registry of interchangeable backends.
\item An optimization ladder of eight rungs (NumPy, eager PyTorch with roll or gather
streaming, \texttt{torch.compile}, and fused kernels in Triton, Warp, Numba-CUDA and CUDA C)
plus memory-traffic variants: AA-pattern in-place streaming~\citep{bailey2009} and 16-bit
storage of shifted populations~\citep{lehmann2022}. Every rung reproduces the reference to
round-off.
\item A kernel generator that prints the same fused update for Triton, Warp and Numba, so
that differences between those backends are purely due to the compilers.
\item A roofline analysis against the measured copy bandwidth showing that Python-driven
kernels reach the same hardware efficiency as hand-written CUDA C.
\item Validation against five reference problems, and an account of the defects of the
original code (mirrored axis, specular side walls, wrong lid density).
\end{itemize}

\section{Background and related work}
\paragraph{LBM on GPUs.}
Single-GPU LB implementations have tracked the memory bandwidth of each GPU generation.
T\"olke and Krafczyk~\citep{tolke2008} demonstrated 3D LBM on desktop GPUs with a D3Q13
model. Obrecht et al.~\citep{obrecht2011} reported up to 512 MLUPS (LBGK) and 516 MLUPS
(MRT) for D3Q19 in single precision on a GeForce GTX~295, reaching up to 86\% of the maximum
global-memory throughput. Mawson and Revell~\citep{mawson2014} reached 1036 MLUPS for D3Q19
BGK in single precision on a Tesla K20c. Lehmann et al.~\citep{lehmann2022} report
5493 MLUPS for D3Q19 with FP32 storage on an RTX~3090 and 15\,455 MLUPS on an A100 with
FP32 arithmetic and FP16 storage of shifted populations.\footnote{Numbers from Table~5 and
the conclusions of the arXiv version (2112.08926v2) of~\citep{lehmann2022}.} Because these
results come from different GPUs, we compare fractions of bandwidth, not MLUPS.

\paragraph{Memory layouts and streaming patterns.}
The two-buffer ``pull'' scheme reads $f_i(\mathbf{x}-\mathbf{c}_i)$ and writes $f_i(\mathbf{x})$,
needing two copies of the populations. In-place schemes avoid the second copy: the AA
pattern~\citep{bailey2009} alternates two kernels that read and write the same locations, and
the esoteric twist~\citep{geier2017} achieves the same with a single kernel. Wittmann et
al.~\citep{wittmann2013} compare propagation steps on CPUs and GPUs. Code-generation
frameworks such as lbmpy~\citep{bauer2021lbmpy} (used inside waLBerla~\citep{bauer2021walberla})
derive optimized kernels symbolically; our generator follows the same idea at a much smaller
scale.

\paragraph{Reduced precision.}
Lehmann et al.~\citep{lehmann2022} show that FP32 arithmetic with 16-bit storage is accurate
for many flows provided the populations are stored shifted, $f_i-w_i$ (the deviation from the
rest equilibrium), which recovers the bits that would otherwise be spent on the constant
$w_i$. We adopt this ``DDF shifting'' for our fp16/bf16 variants.

\paragraph{Python toolchains.}
PyTorch~\citep{paszke2019} with \texttt{torch.compile}~\citep{ansel2024} lowers tensor
programs to Triton~\citep{tillet2019} kernels. Triton, NVIDIA Warp~\citep{warp2022},
Numba~\citep{lam2015} and CuPy~\citep{okuta2017} let one write GPU kernels from Python;
CuPy compiles CUDA C at run time through NVRTC. Lettuce~\citep{bedrunka2021} and
XLB~\citep{ataei2024} are LB frameworks on PyTorch and JAX respectively.

\section{Method}
\subsection{Lattice Boltzmann equation}
With populations $f_i(\mathbf{x},t)$ on a DdQq lattice (D2Q9 or D3Q19, weights $w_i$, lattice
speed of sound $c_s^2=1/3$), one time step is
\begin{equation}
f_i(\mathbf{x}+\mathbf{c}_i, t+1) = f_i^\star(\mathbf{x},t)
= f_i + \Omega_i(\mathbf{f}) + S_i ,
\end{equation}
with density $\rho=\sum_i f_i$ and velocity $\rho\mathbf{u}=\sum_i\mathbf{c}_i f_i+\mathbf{F}/2$.
The equilibrium is the usual second-order polynomial
$f_i^{\rm eq}=w_i\rho[1+3\,\mathbf{c}_i\!\cdot\!\mathbf{u}+\tfrac92(\mathbf{c}_i\!\cdot\!\mathbf{u})^2-\tfrac32u^2]$.
We implement three collision operators:
BGK~\citep{bhatnagar1954,qian1992}, $\Omega_i=-(f_i-f_i^{\rm eq})/\tau$, with kinematic
viscosity $\nu=(\tau-\tfrac12)/3$;
TRT~\citep{ginzburg2008}, which relaxes the symmetric and antisymmetric parts of
$f_i-f_i^{\rm eq}$ with rates $s_+=1/\tau$ and $s_-$ chosen so that the ``magic'' parameter
$\Lambda=(1/s_+-\tfrac12)(1/s_--\tfrac12)=3/16$;
and MRT~\citep{lallemand2000,dhumieres2002} for D2Q9 in the orthogonal moment basis
$(\rho,e,\varepsilon,j_x,q_x,j_y,q_y,p_{xx},p_{xy})$ of Lallemand and Luo, with shear rates
$1/\tau$, energy-flux rates from the same magic relation, and $s_e=1.64$, $s_\varepsilon=1.54$.
Body forces enter through the Guo source term~\citep{guo2002}
$S_i=(1-\tfrac{s}{2})w_i[3(\mathbf{c}_i-\mathbf{u})\!\cdot\!\mathbf{F}+9(\mathbf{c}_i\!\cdot\!\mathbf{u})(\mathbf{c}_i\!\cdot\!\mathbf{F})]$,
applied per relaxation rate for TRT and MRT.

\subsection{Boundary conditions and conventions}
Arrays are stored structure-of-arrays, $f[i,y,x]$ (or $f[i,z,y,x]$), with $x$ the fastest index,
$y$ pointing up (row~0 is the bottom wall) and the lid on the top wall moving in $+x$. Walls
are placed half-way between the last fluid node and the virtual solid node, so a cavity of
$n$ nodes has side length $n$ and $\mathrm{Re}=U n/\nu$. A population leaving the fluid towards
a wall returns in the opposite direction one step later,
\begin{equation}
f_{\bar\imath}(\mathbf{x},t+1)=f_i^\star(\mathbf{x},t)-6w_i\rho_w\,\mathbf{c}_i\!\cdot\!\mathbf{u}_w ,
\end{equation}
where the moving-wall term is Ladd's momentum correction~\citep{ladd1994} with $\rho_w=1$. At
edges and corners the lid takes precedence. The cylinder benchmark additionally uses a
velocity inlet (the same moving-wall rule with the parabolic profile as $\mathbf{u}_w$) and a
pressure outlet that sets the incoming populations to $f_i^{\rm eq}(\rho=1,\mathbf{u})$ with
the velocity of the neighbouring column.

\subsection{Defects of the original code}
Running the original script through the same interface exposed three defects, all fixed in
the new core. (i)~Its streaming moves $c_y=+1$ populations towards decreasing row index while
its plots use \texttt{origin="lower"}, so the lid that appears at the top is physically the
$y=0$ wall; the flow is the mirror image of the intended one. (ii)~The side walls swap
directions $5\leftrightarrow6$ and $8\leftrightarrow7$, which is a specular (free-slip)
reflection, not bounce-back; no-slip was only enforced indirectly by zeroing $\mathbf{u}$
before the next equilibrium. (iii)~The lid density was computed as
$f_0+f_1+f_3+f_2+f_5+f_6+2f_4$ instead of $f_0+f_1+f_3+2(f_4+f_7+f_8)$. With the
correct mirror mapping, the original code deviates from Ghia's Re\,=\,100 profiles by
$L_\infty=\ghiaOneOrigLinf$ (in units of the lid speed) on a $128^2$ grid, against
$\ghiaOneLinf$ for the corrected solver on the same grid (\cref{sec:validation}).

\section{Implementation: the optimization ladder}
\subsection{One interface, one oracle}
A \texttt{Problem} describes the box (per-axis periodic or walled), $\tau$, collision
operator, lid velocity and force; the general solvers additionally accept solid masks,
an inlet profile and an outlet. Every backend implements \texttt{step(n)},
\texttt{populations()} (the post-collision state in canonical layout, as float64) and
\texttt{set\_populations()}. The NumPy reference evaluates the unfused algorithm in float64
with \texttt{np.roll} streaming and slice-assigned wall rules. The same array-level code runs
on torch tensors, which gives the first PyTorch rungs for free.

\subsection{The rungs}
\begin{description}
\item[NumPy (float64).] The reference: one array operation per population and term.
\item[torch eager, roll.] The reference code on CUDA tensors; about 150 kernel launches per
step, so it is launch-bound at every size.
\item[torch eager, gather.] Streaming and all boundary rules collapse into a single
\texttt{torch.take} with a precomputed index table. The table is produced by running the
reference streaming routine on an array holding each entry's own index, so this rung is
boundary-for-boundary identical to the reference by construction, and it supports arbitrary
geometry (we use it for the cylinder benchmark).
\item[torch.compile (original, roll, gather).] The original script, and the roll and
gather formulations, compiled by Inductor into Triton kernels.
\item[Fused kernels.] One kernel per time step (two for the AA pattern) that pulls the
populations of a node, applies bounce-back, collides and stores, in Triton, Warp,
Numba-CUDA and CUDA C (NVRTC via CuPy). Indices are 32-bit; the wall tests and periodic
wraps reduce to comparisons of the node coordinates with the box edges, so no flag field is
read and the traffic stays at the minimum.
\end{description}

\subsection{A single kernel description for three Python DSLs}
The Triton, Warp and Numba kernels are printed by a small generator
(\texttt{pylbm/codegen.py}) from one description of the update. Each population becomes one
select-then-load statement (the select chooses between the neighbour's slot and the own
opposite slot for wall links, so exactly one load is issued), followed by straight-line
moments, equilibria and collision with constants folded in, and one store per population.
The printers differ only in syntax: Triton operates on a block of \texttt{BLOCK} nodes along
$x$ with masked vector loads, Warp and Numba on one node per thread, and Numba needs explicit
\texttt{float32} literals to avoid silent promotion to fp64. The CUDA C kernel is written by
hand with \texttt{\#pragma unroll} loops over constant tables and specialised by preprocessor
definitions; it does not share code with the generator, which makes it an independent check.

\subsection{Memory-traffic variants}
\paragraph{AA pattern.} The two-buffer scheme holds $2q$ values per node. In the AA
pattern~\citep{bailey2009} an ``even'' kernel reads $f_i(\mathbf{x})$ and writes the
post-collision value into slot $\bar\imath$ of the same node; an ``odd'' kernel reads slot
$\bar\imath$ of node $\mathbf{x}-\mathbf{c}_i$ and writes slot $i$ of node
$\mathbf{x}+\mathbf{c}_i$. Each node reads and writes exactly the locations it owns, so the
update is race free in a single buffer. Wall links need care in the odd step on both sides:
a link whose source is behind a wall reads the node's own slot $i$, and a link whose target
is behind a wall writes back into the node's own slot $\bar\imath$ including the moving-wall
term. Our AA kernels are bit-identical to the two-buffer kernels in fp32
(test \texttt{test\_aa\_pattern\_bitwise\_equal\_to\_two\_buffer}). The traffic per update is
unchanged, but the memory footprint halves.
\paragraph{16-bit storage.} We store $h_i=f_i-w_i$ in fp16 or bf16 and compute in fp32,
following~\citep{lehmann2022}: the conversion costs a few instructions and the traffic drops
to \SI{36}{B} per D2Q9 node.

\subsection{Traffic model and measurement}
\label{sec:model}
We use the minimal traffic model: every update reads and writes each population once,
$2qs$ bytes for storage size $s$ (\SI{72}{B} D2Q9/fp32, \SI{36}{B} D2Q9/16-bit,
\SI{152}{B} D3Q19/fp32, \SI{76}{B} D3Q19/16-bit). Boundary handling adds no traffic in the
fused kernels. Effective bandwidth is MLUPS times this byte count, so it is never
inflated. The ceiling is the bandwidth of a device-to-device copy measured at the start of
each benchmark run with a \SI{1}{GiB} \texttt{torch} \texttt{copy\_} (read plus write
counted): \SI{\copyBW}{GB/s} on the \gpuName{}, against a specification value of
\SI{\specBW}{GB/s}. The resulting fp32 D2Q9 roofline is \ceilFpThirtyTwo{} MLUPS, the
16-bit roofline \ceilSixteen{} MLUPS and the D3Q19 fp32 roofline \ceilThreeD{} MLUPS. The
benchmark problem is the lid-driven cavity (all walls active). Timings follow a fixed
protocol: warm-up including JIT compilation, device synchronisation around every timed
region, five repeats of at least \SI{0.5}{s} each, and the median reported with min/max.

\section{Validation}
\label{sec:validation}
\paragraph{Backend equivalence.} After 100--200 steps on small cavities, channels and
periodic boxes (BGK, TRT, MRT, with and without force, 2D and 3D), every fp32 backend
agrees with the float64 reference to \eqMaxFpThirtyTwo{} in the populations, the fp16
variants to \eqMaxFpSixteen{} and the bf16 variant to \eqMaxBf{} (\eqCount{} backend
variants; \texttt{tests/test\_equivalence.py}).

\paragraph{Lid-driven cavity.}
\Cref{fig:ghia} compares the centreline velocities with the tabulated values of Ghia et
al.~\citep{ghia1982} at Re = 100, 400 and 1000 (lid speed $U=0.1$). Errors are computed at
the 15 interior table points per profile (the Re\,=\,400 $v$ value at $x=0.9063$, a widely
reported misprint, is excluded) and given in units of $U$ in \cref{tab:ghia}. They are at
the level of the differences between Ghia's $129^2$ multigrid solution and modern
high-resolution solutions, and do not decrease further with grid refinement at Re\,=\,100
($\ghiaOneSixtyfourLinf$, $\ghiaOneLinf$, $\ghiaOneTwofiftysixLinf$ on $64^2$, $128^2$,
$256^2$). TRT and MRT agree with BGK at Re\,=\,1000, and the AA-pattern kernel reproduces
the two-buffer result exactly. With shifted fp16 storage the Re\,=\,100 error is unchanged
($L_\infty=\ghiaOneHalLinf$) and the Re\,=\,1000 error grows from \ghiaThouLinf{} to
\ghiaThouHalLinf{}, still about 1\% of the lid speed; the velocity field then stagnates at a
round-off noise floor instead of converging to the fp32 tolerance. bf16, with an 8-bit
significand, is clearly not accurate enough ($L_\infty=\ghiaThouBfLinf$ at Re\,=\,1000).

\begin{table}[htbp]
\centering
\caption{Cavity centreline errors against Ghia et al.\ (1982) in units of the lid speed
($L_2$: RMS over the interior table points of both profiles; $L_\infty$: maximum).}
\label{tab:ghia}
\begin{tabular}{llcc}
\toprule
case & solver & $L_2$ & $L_\infty$\\
\midrule
Re 100, $128^2$ & Triton fused, BGK & \ghiaOneLtwo & \ghiaOneLinf\\
Re 400, $256^2$ & Triton fused, BGK & \ghiaFourLtwo & \ghiaFourLinf\\
Re 1000, $256^2$ & Triton fused, BGK & \ghiaThouLtwo & \ghiaThouLinf\\
Re 1000, $256^2$ & Triton fused, TRT & \ghiaThouTRTLtwo & \ghiaThouTRTLinf\\
Re 1000, $256^2$ & Triton fused, MRT & \ghiaThouMRTLtwo & \ghiaThouMRTLinf\\
Re 1000, $256^2$ & Triton fused AA, BGK & \ghiaThouAALtwo & \ghiaThouAALinf\\
\midrule
Re 100, $128^2$ & Triton, fp16 storage & \ghiaOneHalLtwo & \ghiaOneHalLinf\\
Re 100, $128^2$ & Triton, bf16 storage & \ghiaOneBfLtwo & \ghiaOneBfLinf\\
Re 1000, $256^2$ & Triton, fp16 storage & \ghiaThouHalLtwo & \ghiaThouHalLinf\\
Re 1000, $256^2$ & CUDA C, AA + fp16 & \ghiaThouCudaAAHalLtwo & \ghiaThouCudaAAHalLinf\\
Re 1000, $256^2$ & Triton, bf16 storage & \ghiaThouBfLtwo & \ghiaThouBfLinf\\
\midrule
Re 100, $128^2$ & original code & \ghiaOneOrigLtwo & \ghiaOneOrigLinf\\
\bottomrule
\end{tabular}
\end{table}

\begin{figure}[htbp]
\centering
\fig{ghia}
\caption{Lid-driven cavity centreline profiles (lines, finest grid per Re) against Ghia et
al.\ (1982) (circles). The outlying Re\,=\,400 point at $x=0.9063$ is the misprint excluded
from the error norms.}
\label{fig:ghia}
\end{figure}

\paragraph{Shear-wave decay.} A sinusoidal shear wave $u_x=A\sin(2\pi y/n)$ in a periodic
$\shearN^2$ box decays as $\exp(-\nu k^2t)$. Fitting the decay gives the measured viscosity
(\cref{fig:visc}). For $\tau\in[\shearTauMin,\shearTauMax]$ the largest relative deviation
from $(\tau-\tfrac12)/3$ is \shearMaxErr\%, and \shearMidErr\% for $0.6\le\tau\le1.6$; the
remaining deviation is the known higher-order truncation error that grows towards both ends
of the $\tau$ range.

\begin{figure}[htbp]
\centering
\fig{viscosity}
\caption{Shear-wave decay: measured against theoretical viscosity (left) and relative error
(right).}
\label{fig:visc}
\end{figure}

\paragraph{Poiseuille flow.} A body-force driven channel of width $H$ at fixed $\tau=0.8$
and fixed peak velocity converges to the analytic parabola with order \poisOrder{} for BGK
(relative $L_2$ error \poisBgkFinest{} at $H=64$), and TRT with $\Lambda=3/16$ reproduces the
parabola to round-off (largest error \poisTrtMax), as expected for half-way bounce-back with
the magic parameter (\cref{fig:pois}). The convergence study uses the float64 reference:
in fp32 the tiny force increments ($\sim10^{-6}$ per step) are partly lost to round-off when
added to populations of order 0.1, which limits the $H=32$ error to \poisFpThirtyTwo{}
against \poisFpSixtyFourHThirtyTwo{} in float64.

\begin{figure}[htbp]
\centering
\fig{poiseuille}
\caption{Poiseuille channel: grid convergence (left) and profile (right).}
\label{fig:pois}
\end{figure}

\paragraph{Cylinder in a channel.} For the Sch\"afer--Turek benchmark 2D-2~\citep{schafer1996}
(channel $2.2\times0.41$, cylinder diameter $0.1$ at $(0.2,0.2)$, parabolic inflow with mean
velocity $\bar U$, Re\,=\,$\bar UD/\nu=100$) we use $D=\karmanD$ nodes
($\karmanNx\times\karmanNy$ lattice, $\bar U=0.05$, $\tau=\karmanTau$, TRT), a staircase
cylinder and momentum-exchange forces. The flow settles into periodic vortex shedding
(\cref{fig:karman}) with Strouhal number $\mathrm{St}=fD/\bar U=\karmanSt$, maximum drag
coefficient $c_{D,\max}=\karmanCd$ and maximum lift coefficient $c_{L,\max}=\karmanCl$. The
benchmark's reference intervals are $[0.2950,0.3050]$, $[3.22,3.24]$ and $[0.99,1.01]$; the
drag is over-predicted, as expected for a staircase cylinder with simple bounce-back at this
resolution.

\begin{figure}[htbp]
\centering
\fig{karman}
\caption{Sch\"afer--Turek 2D-2 at Re\,=\,100: vorticity (units of $\bar U/D$) and the lift and
drag coefficients over the measurement window.}
\label{fig:karman}
\end{figure}

\paragraph{3D.} A 2D Taylor--Green vortex embedded in each coordinate plane of a periodic
$\tgvN^3$ D3Q19 box decays with the expected viscosity to within \tgvMaxErr\% (Triton and
CUDA C), and the 3D lid-driven cavity of every 3D backend matches the reference.

\section{Performance}
\Cref{fig:ladder} and \cref{tab:ladder} show the ladder for the lid-driven cavity. The
fused kernels are evaluated at the largest benchmark size; the slow rungs at the sizes where
they finish in reasonable time.

\begin{figure}[htbp]
\centering
\fig{ladder}
\caption{The optimization ladder (median, whiskers min/max). Dashed: fp32 roofline at the
measured copy bandwidth (\SI{\copyBW}{GB/s}); dotted: roofline for 16-bit storage.}
\label{fig:ladder}
\end{figure}

\begin{table}[htbp]
\centering
\small
\caption{Lid-driven cavity throughput on the \gpuName{} (median of 5). B/LU: bytes per
lattice update in the minimal traffic model; \% copy and \% spec: effective bandwidth as a
fraction of the measured copy bandwidth and of the \SI{\specBW}{GB/s} specification.}
\label{tab:ladder}
\begin{tabular}{llrrrrrr}
\toprule
 & rung & size & MLUPS & B/LU & GB/s & \% copy & \% spec\\
\midrule
\ladderRows
\\\bottomrule
\end{tabular}
\end{table}

\paragraph{From the original code to the roofline.}
The original script runs at \origMLUPS{} MLUPS (\origPct\% of copy bandwidth): its
\texttt{torch.compile}d functions are separated by Python-level boundary updates on an
array-of-structures layout, so populations make several round trips through DRAM per step.
Compiling the clean roll formulation lets Inductor fuse the step into a few kernels and
gives
\compileRollMLUPS{} MLUPS (\compileRollPct\%). The fused kernels then reach the
roofline: \tritonFusedMLUPS{} MLUPS for Triton (\tritonFusedPct\%), \warpFusedMLUPS{} for
Warp (\warpFusedPct\%), \numbaFusedMLUPS{} for Numba (\numbaFusedPct\%) and
\cudaFusedMLUPS{} for CUDA C (\cudaFusedPct\%). The best fp32 variant is
\speedupFpThirtyTwo$\times$ faster than the original code, and the best 16-bit variant
\speedupBest$\times$. Values slightly above 100\% of the copy bandwidth are possible because
a plain copy does not saturate DRAM perfectly either; the specification bandwidth is a hard
upper bound that no variant exceeds.

\paragraph{Collision operators are free.} At these intensities the arithmetic is hidden
behind memory traffic: the Triton kernel reaches \tritonTRTPct\% (TRT) and \tritonMRTPct\%
(MRT) of copy bandwidth, against \tritonFusedPct\% for BGK.

\paragraph{Problem size.} \Cref{fig:scaling} shows throughput against lattice size. The
fused kernels reach the plateau from about $10^6$ nodes; below that the per-launch overhead
of the Python runtime (tens of microseconds per kernel launch through Triton's and Warp's
launchers) dominates. D3Q19 reaches \bestThreeDMLUPS{} MLUPS (\bestThreeDPct\% of copy
bandwidth) in fp32 and \bestThreeDSixteenMLUPS{} MLUPS (\bestThreeDSixteenPct\%) with
16-bit storage.

\begin{figure}[htbp]
\centering
\fig{scaling}
\caption{Throughput against lattice size for D2Q9 (left) and D3Q19 (right). Solid:
two-buffer fp32; dashed: AA pattern; dotted: 16-bit storage. Horizontal lines: rooflines.}
\label{fig:scaling}
\end{figure}

\paragraph{Comparison with published codes.} Obrecht et al.~\citep{obrecht2011} report up to
86\% of the maximum global-memory throughput for D3Q19 on a GTX~295, and Lehmann et
al.'s~\citep{lehmann2022} RTX~3090 result corresponds to roughly 89\% of that card's
\SI{936}{GB/s} specification bandwidth (5493 MLUPS $\times$ \SI{152}{B}). Our fused kernels
operate at \bestFpThirtyTwoPct\% of measured copy bandwidth on a laptop GPU, i.e.\ in the same
efficiency range as these hand-optimised CUDA/OpenCL codes, and the AA and 16-bit variants
compose with it as expected.

\section{Discussion}
\paragraph{What Python costs.} At the kernel level nothing: once the update is a single fused
kernel, the toolchain matters little and Triton, Warp, Numba and NVRTC-compiled CUDA C all
land on the roofline. Python's cost is in launch overhead (visible at small sizes) and in the
temptation to write the algorithm as many small array operations, which is what limits the
eager and partially compiled rungs. \texttt{torch.compile} removes most of that cost
automatically for the roll formulation, but not for the gather formulation, where the index
table doubles the traffic.

\paragraph{Threats to validity.}
All measurements come from a single RTX~3070 Laptop GPU. Laptop GPUs adjust clocks with
temperature and power; we record the SM clock and temperature with every run, but sustained
runs may throttle. During development the GPU was shared with other jobs; publishable numbers
come from a serial run on an otherwise idle GPU. The copy-bandwidth ceiling is itself a
measurement (\SI{\copyBW}{GB/s}, \SI{1}{GiB} buffers) and a kernel with a better access
pattern can exceed it slightly; we therefore also report the fraction of the specification
bandwidth. CUDA C is compiled with NVRTC through CuPy, not \texttt{nvcc}; both use the same
NVVM back end, but default flags differ. The minimal traffic model ignores the (cached)
re-reads of neighbouring rows and the boundary branches, so it is a lower bound on actual
DRAM traffic. The benchmark is a lid-driven cavity: problems with complex geometry need a
flag field or index lists and move more data.

\paragraph{Limitations and future work.}
The fused kernels handle box geometries (walls, periodicity, moving lid, body force);
the cylinder benchmark uses the gather backend. MRT is implemented for D2Q9 only. Only one
GPU was available, so multi-GPU domain decomposition with halo exchange is left for future
work; we make no scaling claims.

\section{Conclusion}
A lattice-Boltzmann solver written and driven from Python reaches the memory roofline of a
laptop GPU: fused kernels written in Triton, Warp, Numba or CUDA C all sustain close to the
measured copy bandwidth, and 16-bit storage of shifted populations nearly doubles the update
rate at a moderate accuracy cost (fp16) on the cavity benchmark. The path from a few hundred to
several thousand MLUPS was mostly about data movement, not about the language.

\section*{Code availability}
Code, benchmark data and the scripts that generate every number and figure in this paper are
available at \url{https://github.com/walidhemati/pylbm} under the MIT license.
\texttt{bench/run.py --full} reproduces the performance data, \texttt{pylbm validate} the
validation data, and \texttt{figures/make\_figures.py} the figures and
\texttt{paper/results.tex}.

{\small
\setlength{\bibsep}{1.5pt}
\bibliographystyle{unsrtnat}
\bibliography{references}
}

\end{document}
